% PDF page 64
\source{gilbarg-trudinger:page-0064}

that the function $w$ will belong to $C^\infty(\overline{\Omega})$. If, on the other hand, $f$ is merely assumed continuous, the Newtonian potential $w$ is not necessarily twice differentiable. It turns out that a convenient class of functions $f$ to work with is the class of Hölder continuous functions which we introduce now.

\begin{definition}\label{gt:def-holder-pointwise}\dcref{(4.3)}\group{ch04-poisson-newtonian-potential}\source{gilbarg-trudinger:page-0064}
Let $x_0$ be a point in $\mathbb{R}^n$ and $f$ a function defined on a bounded set $D$ containing $x_0$. If $0<\alpha<1$, we say that $f$ is Hölder continuous with exponent $\alpha$ at $x_0$ if the quantity

\[
\tag{4.3}
[f]_{\alpha;x_0}=\sup_D \frac{|f(x)-f(x_0)|}{|x-x_0|^\alpha}
\]

is finite. We call $[f]_{\alpha;x_0}$ the $\alpha$-Hölder coefficient of $f$ at $x_0$ with respect to $D$. Clearly if $f$ is Hölder continuous at $x_0$, then $f$ is continuous at $x_0$. When (4.3) is finite for $\alpha=1$, $f$ is said to be Lipschitz continuous at $x_0$.
\end{definition}

\emph{Example.} The function $f$ on $B_1(0)$ given by $f(x)=|x|^\beta$, $0<\beta<1$, is Hölder continuous with exponent $\beta$ at $x=0$, and is Lipschitz continuous when $\beta=1$.

\begin{definition}\label{gt:def-holder-uniform-local}\dcref{(4.4)}\group{ch04-poisson-newtonian-potential}\source{gilbarg-trudinger:page-0064}
The notion of Hölder continuity is readily extended to the whole of $D$ (not necessarily bounded). We call $f$ uniformly Hölder continuous with exponent $\alpha$ in $D$ if the quantity

\[
\tag{4.4}
[f]_{\alpha;D}=\sup_{\substack{x,y\in D\\x\neq y}}\frac{|f(x)-f(y)|}{|x-y|^\alpha}, \qquad 0<\alpha\le 1,
\]

is finite; and \emph{locally Hölder continuous} with exponent $\alpha$ in $D$ if $f$ is uniformly Hölder continuous with exponent $\alpha$ on compact subsets of $D$. These two concepts obviously coincide when $D$ is compact. Furthermore note that local Hölder continuity is a stronger property than pointwise Hölder continuity in compact subsets. A locally Hölder continuous function will be pointwise Hölder continuous in $D$ provided it is also bounded in $D$.
\end{definition}

Hölder continuity proves to be a quantitative measure of continuity that is especially well suited to the study of partial differential equations. In a certain sense, it may also be viewed as a fractional differentiability. This suggests a natural widening of the well known spaces of differentiable functions. Let $\Omega$ be an open set in $\mathbb{R}^n$ and $k a$ non-negative integer. The Hölder spaces $C^{k,\alpha}(\overline{\Omega})$ ($C^{k,\alpha}(\Omega)$) are defined as the subspaces of $C^k(\overline{\Omega})$ ($C^k(\Omega)$) consisting of functions whose $k$-th order partial derivatives are uniformly Hölder continuous (locally Hölder continuous) with exponent $\alpha$ in $\Omega$. For simplicity we write

\[
C^{0,\alpha}(\Omega)=C^\alpha(\Omega),\qquad C^{0,\alpha}(\overline{\Omega})=C^\alpha(\overline{\Omega}).
\]

with the understanding that $0<\alpha<1$ whenever this notation is used, unless otherwise stated.

\begin{definition}\label{gt:def-holder-spaces}\dcref{(4.5), (4.6), (4.6$'$), (4.7)}\group{ch04-poisson-newtonian-potential}\source{gilbarg-trudinger:page-0065}
Also, by setting

\[
C^{k,0}(\Omega)=C^k(\Omega),\qquad C^{k,0}(\overline{\Omega})=C^k(\overline{\Omega}).
\]

% PDF page 65
\source{gilbarg-trudinger:page-0065}

we may include the $\C^k(\Omega)$ ($\C^k(\bar{\Omega})$) spaces among the $\C^{k,\alpha}(\Omega)$ ($\C^{k,\alpha}(\bar{\Omega})$) spaces for
$0 \le \alpha \le 1$. We also designate by $\C^{k,\alpha}_0(\Omega)$ the space of functions on $\C^{k,\alpha}(\Omega)$ having
compact support in $\Omega$.

\noindent Let us set

\begin{equation}
[u]_{k,0;\Omega} = |D^k u|_{0;\Omega} = \sup_{|\beta| = k}\sup_{\Omega} |D^\beta u|, \qquad k = 0, 1, 2, \ldots
\tag{4.5}
\end{equation}

\begin{equation}
[u]_{k,\alpha;\Omega} = [D^k u]_{\alpha;\Omega} = \sup_{|\beta| = k} [D^\beta u]_{\alpha;\Omega}.
\end{equation}

\noindent With these seminorms, we can define the related norms

\begin{equation}
\|u\|_{\C^k(\bar{\Omega})} = |u|_{k;\Omega} = |u|_{k,0;\Omega} = \sum_{j=0}^k [u]_{j,0;\Omega} = \sum_{j=0}^k |D^j u|_{0;\Omega},
\tag{4.6}
\end{equation}

\begin{equation}
\|u\|_{\C^{k,\alpha}(\bar{\Omega})} = |u|_{k,\alpha;\Omega} = |u|_{k;\Omega} + [u]_{k,\alpha;\Omega}
= |u|_{k;\Omega} + [D^k u]_{\alpha;\Omega},
\end{equation}

\noindent on the spaces $\C^k(\bar{\Omega})$, $\C^{k,\alpha}(\bar{\Omega})$, respectively. It is sometimes useful, especially in
this chapter, to introduce non-dimensional norms on $\C^k(\bar{\Omega})$, $\C^{k,\alpha}(\bar{\Omega})$. If $\Omega$ is
bounded, with $d = \diam \Omega$, we set

\begin{equation}
\|u\|'_{\C^k(\bar{\Omega})} = |u|'_{k;\Omega} = \sum_{j=0}^k d^j [u]_{j,0;\Omega} = \sum_{j=0}^k d^j |D^j u|_{0;\Omega};
\tag{4.6$'$}
\end{equation}

\begin{equation}
\|u\|'_{\C^{k,\alpha}(\bar{\Omega})} = |u|'_{k,\alpha;\Omega} = |u|'_{k;\Omega} + d^{k+\alpha}[u]_{k,\alpha;\Omega}
= |u|'_{k;\Omega} + d^{k+\alpha}[D^k u]_{\alpha;\Omega}.
\end{equation}

\noindent The spaces $\C^k(\bar{\Omega})$, $\C^{k,\alpha}(\bar{\Omega})$, equipped with their respective norms, are Banach
spaces; (see Chapter 5).

\noindent We note here that the product of Hölder continuous functions is again Hölder
continuous. In fact if $u \in \C^\alpha(\bar{\Omega})$, $v \in \C^\beta(\bar{\Omega})$, we have $uv \in \C^\gamma(\bar{\Omega})$ where $\gamma = \min(\alpha,\beta)$,
and

\begin{equation}
\|uv\|_{\C^\gamma(\bar{\Omega})} \le \max\bigl(1, d^{\alpha+\beta-2\gamma}\bigr)\|u\|_{\C^\alpha(\bar{\Omega})}\|v\|_{\C^\beta(\bar{\Omega})};
\tag{4.7}
\end{equation}

\begin{equation}
\|uv\|'_{\C^\gamma(\bar{\Omega})} \le \|u\|'_{\C^\alpha(\bar{\Omega})}\|v\|'_{\C^\beta(\bar{\Omega})}.
\end{equation}

\noindent For the domains $\Omega$ of interest in this work the inclusion relation $\C^{k',\alpha'}(\bar{\Omega}) \subset$
$\C^{k,\alpha}(\bar{\Omega})$ will hold whenever $k+\alpha < k' + \alpha'$. It should be noted, however, that such a
relation will not be true in general. For example, consider the cusped domain

\[
\Omega = \{(x,y) \in \mathbb{R}^2 \mid y < |x|^{1/2}, \ x^2 + y^2 < 1\};
\]

\noindent and for some $\beta$, $1 < \beta < 2$, let $u(x,y) = (\sgn x)\, y^\beta$ if $y > 0$, $= 0$ if $y \le 0$. Clearly
$u \in \C^1(\bar{\Omega})$. However, if $1 > \alpha > \beta/2$, it is easily seen that $u \notin \C^\alpha(\bar{\Omega})$, and hence
$\C^1(\bar{\Omega}) \not\subset \C^\alpha(\bar{\Omega})$.
\end{definition}

% PDF page 66
\source{gilbarg-trudinger:page-0066}

\begin{center}
\textbf{4.2. The Dirichlet Problem for Poisson's Equation}
\end{center}

\noindent We show that if $f$ is bounded and Hölder continuous in the bounded domain $\Omega$, the classical Dirichlet problem for Poisson's equation may be solved under the same boundary conditions for which Laplace's equation is solvable (Theorem 2.14). First we require some differentiability results for the Newtonian potential in bounded domains.

\medskip

\noindent In the following the $D$ operator is always taken with respect to the $x$ variable.

\medskip

\begin{lemma}\label{gt:lem-newtonian-c1}\dcref{Lemma 4.1, (4.8)}\group{ch04-poisson-newtonian-potential}\source{gilbarg-trudinger:page-0066}
\textit{Let $f$ be bounded and integrable in $\Omega$, and let $w$ be the Newtonian potential of $f$. Then $w \in C^1(\R^n)$ and for any $x \in \Omega$,}

\begin{equation}
\tag{4.8}
D_i w(x)=\int_{\Omega} D_i \Gamma(x-y) f(y)\,dy, \qquad i=1,\ldots,n.
\end{equation}

\noindent \textit{Proof.} By virtue of the estimate (2.14) for $D\Gamma$, the function

\begin{equation*}
v(x)=\int_{\Omega} D_i \Gamma(x-y) f(y)\,dy
\end{equation*}

\noindent is well defined. To show that $v=D_i w$, we fix a function $\eta$ in $C^1(\R)$ satisfying $0 \le \eta \le 1$, $0 \le \eta' \le 2$, $\eta(t)=0$ for $t \le 1$, $\eta(t)=1$ for $t \ge 2$ and define for $\varepsilon>0$,

\begin{equation*}
w_{\varepsilon}(x)=\int_{\Omega}\Gamma \eta_{\varepsilon} f(y)\,dy,\qquad \Gamma=\Gamma(x-y),\qquad \eta_{\varepsilon}=\eta(|x-y|/\varepsilon).
\end{equation*}

\noindent Clearly, $w_{\varepsilon} \in C^1(\R^n)$ and

\begin{equation*}
v(x)-D_i w_{\varepsilon}(x)=\int_{|x-y|\le 2\varepsilon} D_i\{(1-\eta_{\varepsilon})\Gamma\}f(y)\,dy
\end{equation*}

\noindent so that

\begin{equation*}
|v(x)-D_i w_{\varepsilon}(x)|\le \sup |f| \int_{|x-y|\le 2\varepsilon}\left(|D_i\Gamma|+\frac{2}{\varepsilon}|\Gamma|\right)\,dy
\end{equation*}

\begin{equation*}
\le \sup |f|\begin{cases}
\dfrac{2n\varepsilon}{n-2} & \text{for } n>2\\[0.8ex]
4\varepsilon(1+|\log 2\varepsilon|) & \text{for } n=2.
\end{cases}
\end{equation*}

\noindent Consequently, $w_{\varepsilon}$ and $D_i w_{\varepsilon}$ converge uniformly in compact subsets of $\R^n$ to $w$ and $v$ respectively as $\varepsilon\to 0$. Hence, $w \in C^1(\R^n)$ and $D_i w = v$. \hfill $\square$
\end{lemma}

% PDF page 67
\source{gilbarg-trudinger:page-0067}

\begin{lemma}\label{gt:lem-newtonian-c2}\dcref{Lemma 4.2, (4.9)}\group{ch04-poisson-newtonian-potential}\source{gilbarg-trudinger:page-0067, gilbarg-trudinger:page-0068}
Let $f$ be bounded and locally H\"older continuous (with exponent $\alpha \le 1$) in $\Omega$, and let $w$ be the Newtonian potential of $f$. Then $w \in C^2(\Omega)$, $\Delta w = f$ in $\Omega$, and for any $x \in \Omega$,
\begin{equation}\tag{4.9}
D_{ij}w(x)=\int_{\Omega_0} D_{ij}\Gamma(x-y)(f(y)-f(x))\,dy
-f(x)\int_{\partial\Omega_0} D_i\Gamma(x-y)\nu_j(y)\,ds_y,\qquad i,j=1,\ldots,n;
\end{equation}
here $\Omega_0$ is any domain containing $\Omega$ for which the divergence theorem holds and $f$ is extended to vanish outside $\Omega$.

\begin{proof}
By virtue of the estimate (2.14) for $D^2\Gamma$ and the pointwise H\"older continuity of $f$ in $\Omega$ the function
\[
u(x)=\int_{\Omega_0} D_{ij}\Gamma(x-y)(f(y)-f(x))\,dy - f(x)\int_{\partial\Omega_0} D_i\Gamma(x-y)\nu_j(y)\,ds_y
\]
is well-defined. Let $v=D_i w$, and define for $\varepsilon>0$
\[
v_\varepsilon(x)=\int_{\Omega} D_i\Gamma_{\eta_\varepsilon}(y)f(y)\,dy,
\]
where $\eta_\varepsilon$ is the function introduced in the preceding lemma. Clearly, $v_\varepsilon\in C^1(\Omega)$, and differentiating, we obtain
\begin{align*}
D_j v_\varepsilon(x) &= \int_{\Omega} D_j(D_i\Gamma_{\eta_\varepsilon})f(y)\,dy \\
&= \int_{\Omega} D_j(D_i\Gamma_{\eta_\varepsilon})(f(y)-f(x))\,dy \\
&\quad + f(x)\int_{\Omega_0} D_j(D_i\Gamma_{\eta_\varepsilon})\,dy \\
&= \int_{\Omega_0} D_j(D_i\Gamma_{\eta_\varepsilon})(f(y)-f(x))\,dy \\
&\quad - f(x)\int_{\partial\Omega_0} D_i\Gamma\,\nu_j(y)\,ds_y
\end{align*}
\end{proof}
\end{lemma}

% PDF page 68
\source{gilbarg-trudinger:page-0068}



provided $\varepsilon$ is sufficiently small. Hence, by subtraction
\[
\left|u(x)-D_j v_\varepsilon(x)\right|
=
\left|
\int_{|x-y|\leq 2\varepsilon}
D_j\{(1-\eta_\varepsilon)D_i\Gamma\}(f(y)-f(x))\,dy
\right|
\]
\[
\leq [f]_{\alpha;x}
\int_{|x-y|\leq 2\varepsilon}
\left(|D_{ij}\Gamma|+\frac{2}{\varepsilon}|D_i\Gamma|\right)|x-y|^\alpha\,dy
\]
\[
\leq \left(\frac{n}{\alpha}+4\right)[f]_{\alpha;x}(2\varepsilon)^\alpha
\]
provided $2\varepsilon<\operatorname{dist}(x,\partial\Omega)$. Consequently $D_j v_\varepsilon$ converges to $u$ uniformly on compact
subsets of $\Omega$ as $\varepsilon\to 0$, and since $v_\varepsilon$ converges uniformly to $v=D_j w$ in $\Omega$, we obtain
$w\in C^2(\Omega)$ and $u=D_{ij}w$. Finally, setting $\Omega_0=B_R(x)$ in (4.9), we have for sufficiently
large $R$,
\[
\Delta w(x)=\frac{1}{n\omega_n R^{n-1}}\,f(x)\int_{|x-y|=R} v_i(y)v_i(y)\,ds_y=f(x).
\]

This completes the proof of Lemma 4.2. $\square$

From Lemmas 4.1, 4.2 and Theorem 2.14 we can now conclude:

\begin{theorem}\label{gt:thm-dirichlet-poisson}\dcref{Theorem 4.3, (4.10)}\group{ch04-poisson-newtonian-potential}\source{gilbarg-trudinger:page-0068}
\textit{Let $\Omega$ be a bounded domain and suppose that each point of $\partial\Omega$ is
regular (with respect to the Laplacian). Then if $f$ is a bounded, locally H\"older con-
tinuous function in $\Omega$, the classical Dirichlet problem: $\Delta u=f$ in $\Omega$, $u=\varphi$ on $\partial\Omega$, is
uniquely solvable for any continuous boundary values $\varphi$.}

\textit{Proof.} We define $w$ to be the Newtonian potential of $f$ and set $v=u-w$. Then
the problem $\Delta u=f$ in $\Omega$, $u=\varphi$ on $\partial\Omega$ is equivalent to the problem $\Delta v=0$ in $\Omega$,
$v=\varphi-w$ on $\partial\Omega$, and its unique solvability follows by Theorem 2.14. $\square$

In the case where $\Omega$ is a ball, $B=B_R(0)$ say, Theorem 4.3 follows from the
Poisson integral formula (Theorem 2.6) and Lemmas 4.1, 4.2. Moreover we have
the explicit formula for the solution:

\begin{equation}
\tag{4.10}
u(x)=\int_{\partial B} K(x,y)\varphi(y)\,ds_y+\int_B G(x,y)f(y)\,dy
\end{equation}
where $K$ is the Poisson kernel (2.29) and $G$ is the Green's function (2.23).
\end{theorem}

\section*{4.3. H\"older Estimates for the Second Derivatives}

The following lemma provides the basic estimate in the subsequent development of
the theory.


% PDF page 69
\source{gilbarg-trudinger:page-0069}



\section*{4.3. H\"older Estimates for the Second Derivatives}

\begin{lemma}\label{gt:lem-newtonian-holder-ball}\dcref{Lemma 4.4, (4.11), (4.12), (4.13)}\group{ch04-poisson-newtonian-potential}\source{gilbarg-trudinger:page-0069, gilbarg-trudinger:page-0070, gilbarg-trudinger:page-0071}
\textit{Let $B_1 = B_R(x_0)$, $B_2 = B_{2R}(x_0)$ be concentric balls in $\mathbb{R}^n$. Suppose
$f \in C^\alpha(\overline{B_2})$, $0<\alpha<1$, and let $w$ be the Newtonian potential of $f$ in $B_2$. Then $w \in$
$C^{2,\alpha}(\overline{B_1})$ and}

\begin{equation}
\tag{4.11}
\left|D^2 w\right|_{0,\alpha;B_1} \leq C\left|f\right|_{0,\alpha;B_2},
\end{equation}

\textit{i.e.,}
\[
\left|D^2 w\right|_{0;B_1} + R^\alpha \left[D^2 w\right]_{\alpha;B_1}
\leq C\left(\left|f\right|_{0;B_2} + R^\alpha \left[f\right]_{\alpha;B_2}\right)
\]
\textit{where $C = C(n,\alpha)$.}

\textit{Proof.} \textit{For any $x \in B_1$, we have by formula (4.9),}

\[
D_{ij}w(x)=\int_{B_2} D_{ij}\Gamma(x-y)(f(y)-f(x))\,dy-f(x)\int_{\partial B_2} D_i\Gamma(x-y)v_j(y)\,ds_y
\]

\textit{so that by (2.14)}

\begin{equation}
\tag{4.12}
\left|D_{ij}w(x)\right|
\leq \frac{|f(x)|}{n\omega_n}R^{1-n}\int_{\partial B_2}ds_y
+ \frac{[f]_{\alpha;x}}{\omega_n}\int_{B_2}|x-y|^{\alpha-n}\,dy
\end{equation}

\[
\leq 2^{\,n-1}|f(x)| + \frac{n}{\alpha}(3R)^\alpha [f]_{\alpha;x}
\]

\[
\leq C_1\bigl(|f(x)| + R^\alpha [f]_{\alpha;x}\bigr)
\]

\textit{where $C_1 = C_1(n,\alpha)$.}
\par
\noindent Next, for any other point $\bar{x} \in B_1$ we have again by formula (4.9),

\[
D_{ij}w(\bar{x})=\int_{B_2} D_{ij}\Gamma(\bar{x}-y)(f(y)-f(\bar{x}))\,dy
\]

\[
-f(\bar{x})\int_{\partial B_2} D_i\Gamma(\bar{x}-y)v_j(y)\,ds_y.
\]

\noindent Writing $\delta = |x-\bar{x}|$, $\xi = \frac{1}{2}(x+\bar{x})$, we consequently obtain by subtraction

\[
D_{ij}w(\bar{x})-D_{ij}w(x)=f(x)I_1+\bigl(f(x)-f(\bar{x})\bigr)I_2+I_3+I_4
\]

\[
+\bigl(f(x)-f(\bar{x})\bigr)I_5+I_6,
\]


% PDF page 70
\source{gilbarg-trudinger:page-0070}

where the integrals $I_1$, $I_2$, $I_3$, $I_4$, $I_5$ and $I_6$ are given by

\begin{equation}
I_1=\int_{\partial B_2}\bigl(D_i\Gamma(x-y)-D_i\Gamma(\bar{x}-y)\bigr)v_j(y)\,ds_y
\end{equation}

\begin{equation}
I_2=\int_{\partial B_2}D_i\Gamma(\bar{x}-y)v_j(y)\,ds_y
\end{equation}

\begin{equation}
I_3=\int_{B_\delta(\xi)}D_{ij}\Gamma(x-y)\bigl(f(x)-f(y)\bigr)\,dy
\end{equation}

\begin{equation}
I_4=\int_{B_\delta(\xi)}D_{ij}\Gamma(\bar{x}-y)\bigl(f(y)-f(\bar{x})\bigr)\,dy
\end{equation}

\begin{equation}
I_5=\int_{B_2-B_\delta(\xi)}D_{ij}\Gamma(x-y)\,dy
\end{equation}

\begin{equation}
I_6=\int_{B_2-B_\delta(\xi)}\bigl(D_{ij}\Gamma(x-y)-D_{ij}\Gamma(\bar{x}-y)\bigr)\bigl(f(\bar{x})-f(y)\bigr)\,dy.
\end{equation}

The estimation of these integrals can be achieved as follows:

\begin{align*}
|I_1|&\le |x-\bar{x}|\int_{\partial B_2}\bigl|DD_i\Gamma(\hat{x}-y)\bigr|\,ds_y
\qquad \text{for some point $\hat{x}$ between $x$ and $\bar{x}$,}\\
&\le \frac{n^2 2^{\,n-1}|x-\bar{x}|}{R}, \qquad \text{since $|\hat{x}-y|\ge R$ for $y\in \partial B_2$,}\\
&\le n^2 2^{\,n-\alpha}\left(\frac{\delta}{R}\right)^\alpha, \qquad \text{since $\delta=|x-\bar{x}|<2R$.}
\end{align*}

\begin{equation}
|I_2|\le \frac{1}{n\omega_n}R^{1-n}\int_{\partial B_2}ds_y=2^{n-1}.
\end{equation}

\begin{align*}
|I_3|&\le \int_{B_\delta(\xi)}\bigl|D_{ij}\Gamma(x-y)\bigr|\bigl|f(x)-f(y)\bigr|\,dy\\
&\le \frac{1}{\omega_n}[f]_{x;x}\int_{B_{3\delta/2}(x)}|x-y|^{1-n}\,dy\\
&= \frac{n}{\alpha}\left(\frac{3\delta}{2}\right)^\alpha [f]_{x;x}
\end{align*}

\begin{equation}
|I_4|\le \frac{n}{\alpha}\left(\frac{3\delta}{2}\right)^\alpha [f]_{x;x}, \qquad \text{as in the estimation of $I_3$.}
\end{equation}

% PDF page 71
\source{gilbarg-trudinger:page-0071}

\setcounter{section}{3}
\noindent\textbf{4.3. \ Hölder Estimates for the Second Derivatives}\hfill 59

\medskip

\noindent\textbf{Integration by parts gives}

\[
|I_5|=\left|\int_{\partial(B_2-B_\delta(\xi))} D_i\Gamma(x-y)v_j(y)\,ds_y\right|
\]

\[
\le \left|\int_{\partial B_2} D_i\Gamma(x-y)v_j(y)\,ds_y\right|
+\left|\int_{\partial B_\delta(\xi)} D_i\Gamma(x-y)v_j(y)\,ds_y\right|
\]

\[
\le 2^{n-1}+\frac{1}{n\omega_n}\left(\frac{\delta}{2}\right)^{1-n}\int_{\partial B_\delta(\xi)}ds_y=2^n.
\]

\[
|I_6|\le |x-\bar{x}|\int_{B_2-B_\delta(\xi)} |DD_{ij}\Gamma(\hat{x}-y)|\,|f(\bar{x})-f(y)|\,dy
\]

\[
\text{for some }\hat{x}\text{ between }x\text{ and }\bar{x},
\]

\[
\le c\delta\int_{|y-\xi|\ge \delta}\frac{|f(\bar{x})-f(y)|}{|\hat{x}-y|^{n+1}}\,dy,\qquad
c=n(n+5)/\omega_n
\]

\[
\le c\delta [f]_{\alpha;x}\int_{|y-\xi|\ge \delta}\frac{|\bar{x}-y|^\alpha}{|\hat{x}-y|^{n+1}}\,dy,
\]

\[
\le c\left(\frac32\right)^\alpha 2^{n+1}\delta [f]_{\alpha;x}\int_{|y-\xi|\ge \delta}|\xi-y|^{\alpha-n-1}\,dy
\]

\[
\text{since }|\bar{x}-y|\le \frac32|\xi-y|\le 3|\hat{x}-y|,
\]

\[
\le \frac{c'}{1-\alpha}2^{n+1}\left(\frac32\right)^\alpha \delta^\alpha [f]_{\alpha;x},\qquad
c'=n^2(n+5).
\]

\medskip

\noindent\textbf{Collecting terms, we thus have}

\[
\tag{4.13}
|D_{ij}w(\bar{x})-D_{ij}w(x)|\le C_2\bigl(R^{-\alpha}|f(x)|+[f]_{\alpha;x}+[f]_{\alpha;x}\bigr)|x-\bar{x}|^\alpha,
\]

\noindent where $C_2$ is a constant depending only on $n$ and $\alpha$. The required estimate then follows by combining (4.12) and (4.13). \hfill $\square$

\medskip

\noindent\textit{Remark.} If $\Omega_1$, $\Omega_2$ are domains such that $\Omega_1\subset B_1$, $\Omega_2\supset B_2$, and $f\in C^\alpha(\overline{\Omega_2})$, and if $w$ is the Newtonian potential of $f$ over $\Omega_2$, then evidently Lemma 4.4 remains true with $\Omega_1$, $\Omega_2$ replacing $B_1$, $B_2$, respectively, in (4.11); that is,

\[
|D^2w|_{0,\alpha;\Omega_1}\le C|f|_{0,\alpha;\Omega_2}.
\]

\medskip

\noindent Hölder estimates for solutions of Poisson's equation follow immediately from Lemma 4.4.

\end{lemma}


% PDF page 72
\source{gilbarg-trudinger:page-0072}


\setcounter{section}{3}
\noindent\textbf{4. Poisson's Equation and the Newtonian Potential}\hfill 60

\medskip

\begin{theorem}\label{gt:thm-poisson-compact-support}\dcref{Theorem 4.5, (4.14), (4.15)}\group{ch04-poisson-newtonian-potential}\source{gilbarg-trudinger:page-0072}
\textit{Let $u\in C_0^2(\R^n)$, $f\in C_0^\alpha(\R^n)$, satisfy Poisson's equation $\Delta u=f$ in $\R^n$. Then $u\in C_0^{2,\alpha}(\R^n)$ and, if $B=B_R(x_0)$ is any ball containing the support of $u$, we have}

\[
\tag{4.14}
\left|D^2u\right|_{0,\alpha;B}\le C\left|f\right|_{0,\alpha;B},\qquad C=C(n,\alpha)
\]

\[
\left|u\right|_{1;B}\le CR^2\left|f\right|_{0;B},\qquad C=C(n)
\]

\textit{Proof.} \textit{By virtue of the representation (2.17),}

\[
\tag{4.15}
u(x)=\int_{\R^n}\Gamma(x-y)f(y)\,dy,
\]

\textit{so that the estimates for $Du$ and $D^2u$ follow respectively from Lemmas 4.1 and 4.4 and the fact that $f$ has compact support in $B$. The estimate for $\left|u\right|_{0;B}$ follows at once from that for $Du$.} $\square$

\medskip

\noindent\textit{The restriction that $u$ has compact support can be removed, by various means, in order to achieve the following interior H\"older estimate for solutions of Poisson's equation. (See also Problem 4.4.)}

\end{theorem}

\medskip

\begin{theorem}\label{gt:thm-poisson-interior-holder}\dcref{Theorem 4.6, (4.16)}\group{ch04-poisson-newtonian-potential}\source{gilbarg-trudinger:page-0072, gilbarg-trudinger:page-0073}
\textit{Let $\Omega$ be a domain in $\R^n$ and let $u\in C^2(\Omega)$, $f\in C^\alpha(\Omega)$, satisfy Poisson's equation $\Delta u=f$ in $\Omega$. Then $u\in C^{2,\alpha}(\Omega)$ and for any two concentric balls $B_1=B_R(x_0)$, $B_2=B_{2R}(x_0)\subset\subset \Omega$ we have}

\[
\tag{4.16}
\left|u\right|_{2,\alpha;B_1}\le C\bigl(\left|u\right|_{0;B_2}+R^2\left|f\right|_{0,\alpha;B_2}\bigr)
\]

\textit{where $C=C(n,\alpha)$.}

\textit{Proof.} \textit{By either Green's representation (2.16) or Lemma 4.2 we can write for $x\in B_2$, $u(x)=v(x)+w(x)$, where $v$ is harmonic in $B_2$ and $w$ is the Newtonian potential of $f$ in $B_2$. By Theorem 2.10 and Lemmas 4.1 and 4.4, we have}

\[
R\left|Dv\right|_{0;B_1}+R^2\left|D^2w\right|_{0,\alpha;B_1}\le CR^2\left|f\right|_{0,\alpha;B_2}
\]

\[
R\left|Dv\right|_{0;B_1}+R^2\left|D^2v\right|_{0,\alpha;B_1}\le C\left|v\right|_{0;B_2}\le C\bigl(\left|u\right|_{0;B_2}+R^2\left|f\right|_{0;B_2}\bigr).
\]

\textit{The last inequality is immediate from $v=u-w$ when $n>2$. For $n=2$, by writing $u(x_1,x_2,x_3)=u(x_1,x_2)$, we may consider $u$ to be a solution of Poisson's equation in a ball in $\R^3$ and the inequality follows in the same way. The desired estimate for $u$ is obtained by combining these inequalities.} $\square$
\end{theorem}

\medskip

\noindent\textit{An immediate consequence of the interior estimate (4.16) is the equicontinuity on compact subdomains of the second derivatives of any bounded set of solutions of Poisson's equation $\Delta u=f$. Consequently by Arzela's theorem we obtain an}


% PDF page 73
\source{gilbarg-trudinger:page-0073}

\noindent
extension of the compactness result, Theorem 2.11, to solutions of Poisson’s
equation.

\medskip

\begin{corollary}\label{gt:cor-poisson-compactness}\dcref{Corollary 4.7}\group{ch04-poisson-newtonian-potential}\source{gilbarg-trudinger:page-0073}
\emph{Any bounded sequence of solutions of Poisson’s equation $\Delta u=f$
in a domain $\Omega$ with $f\in C^{\alpha}(\Omega)$ contains a subsequence converging uniformly on
compact subdomains to a solution.}
\end{corollary}

\medskip

\noindent
It is sometimes preferable to work with an alternative (but equivalent) formula-
tion of the interior estimate (4.16) in terms of certain interior norms which will
be useful later. For $x, y \in \Omega$, which may be any proper open subset of $\mathbb{R}^n$, let
us write $d_x=\dist(x,\partial\Omega)$, $d_{x,y}=\min(d_x,d_y)$. We define for $u\in C^k(\Omega), C^{k,\alpha}(\Omega)$
the following quantities, which are the analogues of the global seminorms and
norms (4.5), (4.6).

\[
[u]^*_{k,0;\Omega}=[u]^*_{k;\Omega}=\sup_{\substack{x\in\Omega\\ |\beta|=k}} d_x^k|D^\beta u(x)|,\qquad k=0,1,2,\ldots;
\]

\[
\norm{u}^*_{k;\Omega}=\norm{u}^*_{k,0;\Omega}=\sum_{j=0}^k [u]^*_{j;\Omega};
\]

\begin{equation}\tag{4.17}
[u]^*_{k,\alpha;\Omega}=\sup_{\substack{x,y\in\Omega\\ |\beta|=k}} d_{x,y}^{k+\alpha}\frac{|D^\beta u(x)-D^\beta u(y)|}{|x-y|^\alpha},
\qquad 0<\alpha\leq 1;
\end{equation}

\[
\norm{u}^*_{k,\alpha;\Omega}=\norm{u}^*_{k;\Omega}+[u]^*_{k,\alpha;\Omega}.
\]

\noindent
In this notation,

\[
[u]^*_{0;\Omega}=\norm{u}^*_{0;\Omega}=\norm{u}_{0;\Omega}.
\]

\noindent
We note that $\norm{u}^*_{k;\Omega}$ and $\norm{u}^*_{k,\alpha;\Omega}$ are norms on the subspaces of $C^k(\Omega)$ and $C^{k,\alpha}(\Omega)$
respectively for which they are finite. If $\Omega$ is bounded and $d=\diam \Omega$, then obviously
these interior norms and the global norms (4.6) are related by

\begin{equation}\tag{4.17$'$}
\norm{u}^*_{k;\Omega}\leq \max(1,d^{k+\alpha})\norm{u}_{k;\Omega}.
\end{equation}

\noindent
If $\Omega' \subset\subset \Omega$ and $\sigma=\dist(\Omega',\partial\Omega)$, then

\begin{equation}\tag{4.17$''$}
\min(1,\sigma^{k+\alpha})\norm{u}_{k,\alpha;\Omega}\leq \norm{u}^*_{k,\alpha;\Omega}.
\end{equation}

\noindent
It is convenient here to also introduce the quantity

\begin{equation}\tag{4.18}
|f|^{(k)}_{0,\alpha;\Omega}=\sup_{x\in\Omega} d_x^k|f(x)|+\sup_{\substack{x,y\in\Omega}} d_{x,y}^{k+\alpha}\frac{|f(x)-f(y)|}{|x-y|^\alpha}.
\end{equation}

\noindent
This is a special case of certain norms to be defined later.
From Theorem 4.6, we can now derive an interior estimate for arbitrary domains
which will be generalized in very similar form to elliptic equations in Chapter 6.


% PDF page 74
\source{gilbarg-trudinger:page-0074}


\begin{theorem}\label{gt:thm-poisson-interior-weighted}\dcref{Theorem 4.8, (4.19), (4.20)}\group{ch04-poisson-newtonian-potential}\source{gilbarg-trudinger:page-0074}
Let $u \in C^2(\Omega)$, $f \in C^\alpha(\Omega)$ satisfy $\Delta u = f$ in an open set $\Omega$ of $\R^n$. Then
\begin{equation}
|u|_{2,\alpha;\Omega}^* \le C\bigl(|u|_{0;\Omega}+|f|_{0,\alpha;\Omega}^{(2)}\bigr),
\tag{4.19}
\end{equation}
where $C=C(n,\alpha)$.
\end{theorem}

\begin{proof}
If either $|u|_{0,\Omega}$ or $|f|_{0,\alpha;\Omega}^{(2)}$ is infinite, the estimate (4.19) is trivial. Otherwise for $x \in \Omega$, $R=\frac{1}{3}d_x$, $B_1=B_R(x)$, $B_2=B_{2R}(x)$, we have for any first derivative $Du$ and second derivative $D^2u$
\begin{align*}
d_x|Du(x)| + d_x^2|D^2u(x)| &\le (3R)|Du|_{0;B_1} + (3R)^2|D^2u|_{0;B_1}\\
&\le C\bigl(|u|_{0;B_2}+R^2|f|_{0,\alpha;B_2}^{(2)}\bigr) \qquad \text{by (4.16)}\\
&\le C\bigl(|u|_{0;\Omega}+|f|_{0,\alpha;\Omega}^{(2)}\bigr).
\end{align*}
Hence we obtain
\begin{equation}
|u|_{2;\Omega}^* \le C\bigl(|u|_{0;\Omega}+|f|_{0,\alpha;\Omega}^{(2)}\bigr).
\tag{4.20}
\end{equation}

To estimate $[u]^*_{2,\alpha;\Omega}$ we let $x$, $y \in \Omega$ with $d_x \le d_y$. Then
\begin{align*}
d_{x,y}^{2+\alpha}\frac{|D^2u(x)-D^2u(y)|}{|x-y|^\alpha}
&\le (3R)^{2+\alpha}[D^2u]_{\alpha;B_1}\\
&\quad + 3^\alpha(3R)^2\bigl(|D^2u(x)|+|D^2u(y)|\bigr)\\
&\le C\bigl(|u|_{0;B_2}+R^2|f|_{0,\alpha;B_2}^{(2)}\bigr) + 6[u]^*_{2;\Omega} \qquad \text{by (4.16)}\\
&\le C\bigl(|u|_{0;\Omega}+|f|_{0,\alpha;\Omega}^{(2)}\bigr) \qquad \text{by (4.20)}.
\end{align*}
The estimate (4.19) then follows. $\square$
\end{proof}

The preceding result provides bounds in compact subsets for $Du$, $D^2u$ and the Hölder coefficients of $D^2u$ in terms of a bound on the right member of (4.19), and hence it is the basis of compactness results for solutions of Poisson's equation. In particular, Corollary 4.7 is also an immediate consequence of Theorem 4.8, after noting that the latter implies the equicontinuity of solutions and of their first and second derivatives on compact subsets.

By means of the compactness result, Corollary 4.7, we can now derive an existence theorem for Poisson's equation $\Delta u=f$ for unbounded $f$.

\begin{theorem}\label{gt:thm-poisson-ball-unbounded}\dcref{Theorem 4.9, (4.21), (4.22)}\group{ch04-poisson-newtonian-potential}\source{gilbarg-trudinger:page-0074, gilbarg-trudinger:page-0075}
Let $B$ be a ball in $\R^n$ and $f$ a function in $C^\alpha(B)$ for which $\sup_{x\in B} d_x^{2-\beta}|f(x)| \le N < \infty$ for some $\beta$, $0<\beta<1$. Then there exists a unique function $u \in C^2(B)\cap C^0(\bar{B})$ satisfying $\Delta u=f$ in $B$, $u=0$ on $\partial B$. Furthermore, $u$ satisfies an estimate
\begin{equation}
\sup_{x\in B} d_x^{-\beta}|u(x)| \le CN,
\tag{4.21}
\end{equation}
where the constant $C$ depends only on $\beta$.
\end{theorem}

% PDF page 75
\source{gilbarg-trudinger:page-0075}
\pagestyle{empty}

\section*{4.3. \ H\"older Estimates for the Second Derivatives}

\textit{Proof.} \ The estimate (4.21) follows from a simple barrier argument. Namely,
let $B=B_R(x_0)$, $r=|x-x_0|$ and set

\[
w(x)=(R^2-r^2)^\beta.
\]

By direct calculation, we have for $r<R$

\[
\Delta w(x)=-2\beta(R^2-r^2)^{\beta-2}\bigl[\ln(R^2-r^2)+2(1-\beta)r^2\bigr]
\le -4\beta(1-\beta)R^2(R^2-r^2)^{\beta-2}\le -\beta(1-\beta)R^\beta(R-r)^{\beta-2}.
\]

Now suppose that $\Delta u=f$ in $B$, $u=0$ on $\partial B$. Since $d_x=R-r$, we have by hypothesis

\[
|f(x)|\le N\, d_x^{\beta-2}=N(R-r)^{\beta-2}
\le -C_0N\,\Delta w,\ \text{where } C_0=[\beta(1-\beta)R^\beta]^{-1},
\]

so that

\[
\Delta(C_0Nw\pm u)\le 0\ \text{in } B \qquad \text{and} \qquad C_0Nw\pm u=0 \ \text{on } \partial B.
\]

Consequently, by the maximum principle,

\[
(4.22)\qquad |u(x)|\le C_0Nw(x)\le CN d_x^\beta \quad \text{for } x\in B,
\]

which implies (4.21) with the constant $C=2/[\beta(1-\beta)]$.
Finally to show the existence of $u$, we let

\[
f_m=
\begin{cases}
m & \text{if } f\ge m \\
f & \text{if } |f|\le m \\
-m & \text{if } f\le -m
\end{cases}
\]

and let $\{B_k\}$ be a sequence of concentric balls exhausting $B$ such that $|f|\le k$ in
$B_k$. We define $u_m$ by $\Delta u_m=f_m$ in $B$, $u_m=0$ on $\partial B$.
By (4.21), we have

\[
\sup_{x\in B} d_x^{-\beta}|u_m(x)|\le C\sup_{x\in B} d_x^{2-\beta}|f_m(x)|\le CN,
\]

so that the sequence $\{u_m\}$ is uniformly bounded and $\Delta u_m=f$ in $B_k$, for $m\ge k$.
Hence by Corollary 4.7, applied successively to the sequence of balls $B_k$, a subse-
quence of $\{u_m\}$ converges in $B$ to a $C^2(B)$ function $u$ satisfying $\Delta u=f$ in $B$. It
follows that $|u(x)|\le CN d_x^\beta$ and hence $u=0$ on $\partial B$. $\square$

It is easy to show by counterexample that Theorem 4.9 is false if $\beta\le 0$. We
remark that the theorem may be extended to more general domains than balls;
(see Problem 4.6). Also, for arbitrary domains with regular boundary points, the
classical Dirichlet problem for Poisson's equation, $\Delta u=f$, is solvable for unbounded
$f$ satisfying certain integrability conditions; (see Problem 4.3).


% PDF page 76
\source{gilbarg-trudinger:page-0076}

\section*{4.4. Estimates at the Boundary}

Theorem 4.8 will be applied in Chapter 6 to the derivation of interior H\"older
estimates for linear elliptic equations. However in order to establish global
estimates, which are required for the existence theory, we need a version of Theorem
4.8 applicable to the intersection of a domain $\Omega$ and a half-space. Let us first derive
the appropriate extension of the H\"older estimate for the Newtonian potential,
Lemma 4.4. In what follows, $\R^{n}_{+}$ will denote the half-space, $x_n>0$, and $T$ the
hyperplane, $x_n=0$; $B_2=B_{2R}(x_0)$, $B_1=B_R(x_0)$ will be balls with center $x_0\in \R^{n}_{+}$
and we let $B_2^+=B_2\cap \R^{n}_{+}$, $B_1^+=B_1\cap \R^{n}_{+}$.

\medskip

\begin{lemma}\label{gt:lem-newtonian-holder-halfball}\dcref{Lemma 4.10, (4.23)}\group{ch04-poisson-newtonian-potential}\source{gilbarg-trudinger:page-0076}
\textit{Let $f\in C^{\alpha}(\bar{B}_2^+)$, and let $w$ be the Newtonian potential of $f$ in $B_2^+$. Then}
\[
w\in C^{2,\alpha}(\bar{B}_1^+) \textit{ and}
\]
\begin{equation}
\tag{4.23}
\lvert D^2w\rvert'_{0,\alpha;B_1^+}\leq C\lvert f\rvert'_{0,\alpha;B_2^+}
\end{equation}
\[
\textit{where }C=C(n,\alpha).
\]

\textit{Proof.} We assume that $B_2$ intersects $T$ since otherwise the result is already
contained in Lemma 4.4. The representation (4.9) holds for $D_{ij}w$ with $\Omega_0=B_2^+$. If
either $i$ or $j\neq n$, then the portion of the boundary integral
\[
\int_{\partial B_2^+} D_i\Gamma(x-y)v_j(y)\,ds_y
\left(
=\int_{\partial B_2^+} D_j\Gamma(x-y)v_i(y)\,ds_y
\right)
\]
on $T$ vanishes since $v_i$ or $v_j=0$ there. The estimates in Lemma 4.4 for $D_{ij}w$
(i or $j\neq n$) then proceed exactly as before with $B_2$ replaced by $B_2^+$, $B_\delta(\xi)$ replaced by
$B_\delta(\xi)\cap B_2^+$ and $\partial B_2$ replaced by $\partial B_2^+-T$. Finally $D_{nn}w$ can be estimated from the
equation $\Delta w=f$ and the estimates on $D_{kk}w$ for $k=1,\ldots,n-1$. $\square$

\medskip

\end{lemma}

\begin{theorem}\label{gt:thm-poisson-halfball-holder}\dcref{Theorem 4.11, (4.24), (4.25), (4.26), (4.27), (4.28)}\group{ch04-poisson-newtonian-potential}\source{gilbarg-trudinger:page-0076, gilbarg-trudinger:page-0077}
\textit{Let $u\in C^2(B_2^+)\cap C^0(\bar{B}_2^+)$, $f\in C^{\alpha}(\bar{B}_2^+)$, satisfy $\Delta u=f$ in $B_2^+$, $u=0$
on $T$. Then $u\in C^{2,\alpha}(\bar{B}_1^+)$ and we have}
\begin{equation}
\tag{4.24}
\lvert u\rvert'_{2,\alpha;B_1^+}\leq C(\lvert u\rvert_{0;B_2^+}+R^2\lvert f\rvert'_{0,\alpha;B_2^+})
\end{equation}
\[
\textit{where }C=C(n,\alpha).
\]

\textit{Proof.} Let $x'=(x_1,\ldots,x_{n-1})$, $x^*=(x',-x_n)$ and define
\[
f^*(x)=f^*(x',x_n)=
\begin{cases}
f(x',x_n) & \text{if }x_n\geq 0\\
f(x',-x_n) & \text{if }x_n\leq 0.
\end{cases}
\]

We assume that $B_2$ intersects $T$; otherwise Theorem 4.6 implies (4.24). We set
$B_2^-=\{x\in \R^n\mid x^*\in B_2^+\}$ and $D=B_2^+\cup B_2^-\cup (B_2\cap T)$. Then $f^*\in C^{\alpha}(D)$ and
\[
\lvert f^*\rvert'_{0,\alpha;D}\leq 2\lvert f\rvert'_{0,\alpha;B_2^+}.
\]

Now, defining

% PDF page 77
\source{gilbarg-trudinger:page-0077}

\begin{equation}
\tag{4.25}
w(x)=\int_{B_2^+}\bigl[\Gamma(x-y)-\Gamma(x^*-y)\bigr]f(y)\,dy
=\int_{B_2^+}\bigl[\Gamma(x-y)-\Gamma(x-y^*)\bigr]f(y)\,dy,
\end{equation}

we have $w(x',0)=0$ (see Problem 2.3c) and $\Delta w=f$ in $B_2^+$. Noting that

\[
\int_{B_2^+}\Gamma(x-y^*)f(y)\,dy=\int_{B_2^-}\Gamma(x-y)f^*(y)\,dy,
\]

we then obtain

\begin{equation}
w(x)=2\int_{B_2^+}\Gamma(x-y)f(y)\,dy-\int_D\Gamma(x-y)f^*(y)\,dy.
\end{equation}

Letting
\[
w^*(x)=\int_D\Gamma(x-y)f^*(y)\,dy,
\]
we have by the Remark following Lemma

4.4 (in which we set $\Omega_1=B_1^+$, $\Omega_2=D$)

\[
|D^2w^*|_{0,\alpha;B_1^+}\le C|f^*|_{0,\alpha;D}\le 2C|f|_{0,\alpha;B_2^+}.
\]

Combining this with Lemma 4.10, we obtain

\begin{equation}
\tag{4.26}
|D^2w|_{0,\alpha;B_1^+}\le C|f|_{0,\alpha;B_2^+}.
\end{equation}

Now let $v=u-w$. Then $\Delta v=0$ in $B_2^+$ and $v=0$ on $T$. By reflection $v$ may be extended to a harmonic function in $B_2$ (Problem 2.4) and hence the estimate (4.24) follows from the interior derivative estimate for harmonic functions, Theorem 2.10. $\square$

\textit{Remark.} If in addition to the hypotheses of Theorem 4.11, $u$ has compact support in $B_2^+\cup T$, we obtain from (4.26) the simpler estimate (extending (4.14))

\begin{equation}
\tag{4.27}
|D^2u|_{0,\alpha;B_2^+}\le C|f|_{0,\alpha;B_2^+}.
\end{equation}

In this case we have the representation

\begin{equation}
\tag{4.28}
u(x)=w(x)=\int_{B_2^+}\bigl[\Gamma(x-y)-\Gamma(x^*-y)\bigr]f(y)\,dy.
\end{equation}

It will be useful to have an analogue of Theorem 4.8 in which the estimates are valid up to a hyperplane boundary piece. For this purpose we introduce certain partially interior norms and seminorms, analogous to (4.17) and (4.18). Let $\Omega$ be a

% PDF page 78
\source{gilbarg-trudinger:page-0078}

proper open subset of $\R^{n}_{+}$ with open boundary portion $T$ on $x_n=0$. For $x,y\in \Omega$
let us write

\[
 d_x=\operatorname{dist}(x,\partial\Omega-T),\ d_{x,y}=\min(d_x,d_y).
\]

We define the following quantities:

\[
[u]^*_{k,0;\Omega\cup T}=[u]^*_{k;\Omega\cup T}=\sup_{\substack{x\in\Omega\\ |\beta|=k}} d_x^{\,k}|D^{\beta}u(x)|,\qquad k=0,1,2,\dots;
\]

\begin{equation}
\tag{4.29}
|u|^*_{k;\Omega\cup T}=|u|^*_{k;0,\Omega\cup T}=\sum_{j=0}^{k}[u]^*_{j;\Omega\cup T};
\end{equation}

\[
[u]^*_{k,\alpha;\Omega\cup T}=\sup_{\substack{x,y\in\Omega\\ |\beta|=k}} d_{x,y}^{\,k+\alpha}\frac{|D^{\beta}u(x)-D^{\beta}u(y)|}{|x-y|^{\alpha}},\qquad 0<\alpha\leq 1;
\]

\[
|u|^*_{k,\alpha;\Omega\cup T}=|u|^*_{k;\Omega\cup T}+[u]^*_{k,\alpha;\Omega\cup T};
\]

\[
|u|^*_{0,\alpha;\Omega\cup T}=\sup_{x\in\Omega} d_x^{\alpha}|u(x)|+\sup_{x,y\in\Omega} d_{x,y}^{\alpha}\frac{|u(x)-u(y)|}{|x-y|^{\alpha}}.
\]

We can now state:

\end{theorem}

\begin{theorem}\label{gt:thm-poisson-partial-boundary}\dcref{Theorem 4.12, (4.29), (4.30)}\group{ch04-poisson-newtonian-potential}\source{gilbarg-trudinger:page-0078}
\quad \textit{Let $\Omega$ be an open set in $\R^{n}_{+}$ with a boundary portion $T$ on $x_n=0$,}
\textit{and let $u\in C^{2}(\Omega)\cap C^{0}(\Omega\cup T)$, $f\in C^{\alpha}(\Omega\cup T)$ satisfy $\Delta u=f$ in $\Omega$, $u=0$ on $T$.}
\textit{Then}

\begin{equation}
\tag{4.30}
|u|^*_{2,\alpha;\Omega\cup T}\leq C\bigl(|u|_{0;\Omega}+|f|^{(2)}_{0,\alpha;\Omega\cup T}\bigr),
\end{equation}

\textit{where $C=C(n,\alpha)$.}

This result follows from Theorem 4.11 in the same way that Theorem 4.8
follows from Theorem 4.6; the details of proof are therefore omitted.
Theorems 4.11 and 4.12 provide a regularity result for solutions of Poisson's
equation at a hyperplane portion of the boundary. More generally, if $\Omega$ is a
bounded domain, $f\in C^{\alpha}(\bar{\Omega})$, $u\in C^{2}(\Omega)\cap C^{0}(\bar{\Omega})$, $\Delta u=f$ in $\Omega$, and if $\partial\Omega$ and the
boundary values of $u$ are sufficiently smooth, it follows that $u\in C^{2,\alpha}(\bar{\Omega})$. This result,
essentially Kellogg's theorem [KE 1], will be established as a byproduct of our
treatment of linear elliptic equations in Chapter 6. The case when $\Omega$ is a ball
however is directly derivable from Theorem 4.11.

\medskip

\end{theorem}

\begin{theorem}\label{gt:thm-poisson-ball-boundary}\dcref{Theorem 4.13, (4.31), (4.32)}\group{ch04-poisson-newtonian-potential}\source{gilbarg-trudinger:page-0079}
\quad \textit{Let $B$ be a ball in $\R^{n}$ and $u$ and $f$ functions on $\bar{B}$ satisfying $u\in$
$C^{2}(B)\cap C^{0}(\bar{B})$, $f\in C^{\alpha}(\bar{B})$, $\Delta u=f$ in $B$, $u=0$ on $\partial B$. Then $u\in C^{2,\alpha}(\bar{B})$.}

\textit{Proof.} By a translation we may assume $\partial B$ passes through the origin. The inversion
mapping $x\to x^{*}=x/|x|^{2}$ is a bicontinuous, smooth mapping of $\R^{n}-\{0\}$ onto


% PDF page 79
\source{gilbarg-trudinger:page-0079}

itself which maps $B$ onto a half-space, $B^{*}$. Furthermore if $u\in C^{2}(B)\cap C^{0}(\overline{B})$, the Kelvin transform of $u$, defined by

\begin{equation}
\tag{4.31}
v(x)=|x|^{2-n}u\!\left(\frac{x}{|x|^{2}}\right)
\end{equation}

belongs to $C^{2}(B^{*})\cap C^{0}(\overline{B^{*}})$ and satisfies (see Problem 4.7).

\begin{equation}
\tag{4.32}
\Delta_{x^{*}}v(x^{*})=|x^{*}|^{-n-2}\Delta_{x}u(x), \qquad x^{*}\in B^{*},\ x\in B,
\end{equation}
\[
=|x^{*}|^{-n-2}f\!\left(\frac{x^{*}}{|x^{*}|^{2}}\right), \qquad x^{*}\in B^{*}.
\]

Hence Theorem 4.11 is applicable to the Kelvin transform $v$ and since by translation any point of $\partial B$ may be taken for the origin we obtain $u\in C^{2,\alpha}(\overline{B})$. $\square$

\end{theorem}

\begin{corollary}\label{gt:cor-poisson-ball-dirichlet}\dcref{Corollary 4.14}\group{ch04-poisson-newtonian-potential}\source{gilbarg-trudinger:page-0079}
\textit{Let $\varphi\in C^{2,\alpha}(\overline{B})$, $f\in C^{\alpha}(\overline{B})$. Then the Dirichlet problem, $\Delta u=f$ in $B$, $u=\varphi$ on $\partial B$, is uniquely solvable for a function $u\in C^{2,\alpha}(\overline{B})$.}

\textit{Proof.} Writing $v=u-\varphi$, the problem is reduced to the problem $\Delta v=f-\Delta\varphi$ in $B$, $v=0$ on $\partial B$, which is solvable for $v\in C^{2}(B)\cap C^{0}(\overline{B})$ by Theorem 4.3 and consequently for $v\in C^{2,\alpha}(\overline{B})$ by Theorem 4.13. $\square$

As a byproduct of the proof of Theorem 4.13, we see that Lemma 4.4 may be improved in the sense that if $f\in C^{\alpha}(\overline{B})$, its Newtonian potential in $B$ will belong to $C^{2,\alpha}(\overline{B})$.

\end{corollary}

\subsection*{4.5. Hölder Estimates for the First Derivatives}

Poisson’s equation often appears in the form

\begin{equation}
\tag{4.33}
\Delta u=\operatorname{div}f=D_{i}f^{i}, \qquad f=(f^{1},\ldots,f^{n})
\end{equation}

where the density function is a divergence. The corresponding estimates of solutions can be reduced to those of the preceding sections, with certain generalizations that will be useful later.

If $f\in C^{1,\alpha}(\Omega)$, then, obviously, the estimates for the Newtonian potential of $\operatorname{div} f$ and for solutions of (4.33) are the same as before provided $\operatorname{div} f$ replaces $f$ throughout. If $\Omega$ is sufficiently smooth, we have

\[
\int_{\Omega}\Gamma(x-y)\,\operatorname{div}f(y)\,dy
=\int_{\Omega}D\Gamma(x-y)\cdot f(y)\,dy+\int_{\partial\Omega}\Gamma(x-y)\,f\cdot \nu\,dS_{y},
\]

% PDF page 80
\source{gilbarg-trudinger:page-0080}


and, thus, the Newtonian potential of $\operatorname{div} f$ in $\Omega$ is within a harmonic function given by

\begin{equation}
\tag{4.34}
w(x)=D\int_{\Omega}\Gamma(x-y)f(y)\,dy=D_{j}\int_{\Omega}\Gamma(x-y)f^{j}(y)\,dy.
\end{equation}

This expression is identical with the Newtonian potential when $f$ has compact support in $\Omega$. We see that $w$ is still defined when $f$ is only integrable, in which case (4.34) can be taken as the definition of a generalized Newtonian potential of $\operatorname{div} f$ in $\Omega$. If, in addition, $f$ is Hölder continuous, the first derivatives of $w$ in $\Omega$ are (as in Lemma 4.2) given by

\begin{equation}
\tag{4.35}
D_{i}w(x)=\int_{\Omega}D_{ij}\Gamma(x-y)\bigl(f^{j}(y)-f^{j}(x)\bigr)\,dy-f^{j}(x)\int_{\partial\Omega}D_{i}\Gamma(x-y)\nu_{j}\,ds_{y},
\end{equation}

which can be estimated as in Lemma 4.4. Thus we have in the notation of Lemma 4.4

\begin{equation}
\tag{4.36}
\lvert Dw\rvert'_{0,\alpha;B_{1}}\leq C\lvert f\rvert'_{0,\alpha;B_{2}}, \qquad C=C(n,\alpha).
\end{equation}

We can, therefore, assert a $C^{1,\alpha}$ interior estimate:

\begin{theorem}\label{gt:thm-poisson-first-derivative-interior}\dcref{Theorem 4.15, (4.33), (4.34), (4.35), (4.36), (4.37)}\group{ch04-poisson-newtonian-potential}\source{gilbarg-trudinger:page-0080}
\textit{Let $\Omega$ be a domain in $\R^{n}$, and let $u$ satisfy Poisson's equation (4.33), where $f\in C^{\alpha}(\Omega)$, $0<\alpha<1$. Then for any two concentric balls $B_{1}=B_{R}(x_{0}), B_{2}=B_{2R}(x_{0})\Subset \Omega$ we have}

\begin{equation}
\tag{4.37}
\lvert u\rvert'_{1,\alpha;B_{1}}\leq C\bigl(\lvert u\rvert_{0;B_{2}}+R\lvert f\rvert'_{0,\alpha;B_{2}}\bigr), \qquad C=C(n,\alpha).
\end{equation}

The proof is the same as that of Theorem 4.6 if (4.36) replaces (4.11) in the argument.

Corresponding boundary estimates can be derived in a similar way. If $B_{1}$ and $B_{2}$ are as in Lemma 4.10, then the potential (4.34) has first derivatives given by (4.35), with $\Omega=B_{2}^{+}$, and we obtain the estimate

\begin{equation}
\tag{4.38}
\lvert Dw\rvert'_{0,\alpha;B_{1}^{+}}\leq C\lvert f\rvert'_{0,\alpha;B_{2}^{+}}, \qquad C=C(n,\alpha).
\end{equation}

To obtain a $C^{1,\alpha}$ analogue of Theorem 4.11 for solutions of (4.33) vanishing on $x_{n}=0$, we proceed as in that theorem with a method based on reflection. Let

\[
G(x,y)=\Gamma(x-y)-\Gamma(x-y^{*})=\Gamma(x-y)-\Gamma(x^{*}-y)
\]

denote the Green's function of the half-space $\R_{+}^{n}$, and consider

\begin{equation}
\tag{4.39}
v(x)=-\int_{B_{2}^{+}}D_{y}G(x,y)\cdot f(y)\,dy
\qquad D_{y}=(D_{y_{1}},\ldots,D_{y_{n}})
\end{equation}

\[
=\int_{B_{2}^{+}}D\Gamma(x-y)\cdot f(y)\,dy+\int_{B_{2}^{+}}D_{y}\Gamma(x^{*}-y)\cdot f(y)\,dy.
\]

% PDF page 81
\source{gilbarg-trudinger:page-0081}

\subsection*{4.5. Hölder Estimates for the First Derivatives}

For each $i = 1,\ldots,n$, let $v_i$ denote the component of $v$ given by

\[
v_i(x)=\int_{B_2^+} D_i\Gamma(x-y)f^i(y)\,dy+\int_{B_2^+} D_{y_i}\Gamma(x^*-y)f^i(y)\,dy.
\]

We see that $v$ and $v_i$ vanish on $T=B_2\cap\{x_n=0\}$. Now suppose $f\in C^\alpha(\overline{B_2^+})$, and let $f$ be extended by even reflection in $x_n=0$, and denote the extended function again by $f$. Then, for $i=1,\ldots,n-1$, we have as in the proof of Theorem 4.11,

\begin{equation}
\tag{4.40}
v_i(x)=D_i\left[2\int_{B_2^+}\Gamma(x-y)f^i(y)\,dy-\int_{B_2^+\cup B_2^-}\Gamma(x-y)f^i(y)\,dy\right].
\end{equation}

And when $i=n$, since

\[
\int_{B_2^+} D_{y_n}\Gamma(x^*-y)f^n(y)\,dy=\int_{B_2^-} D_n\Gamma(x-y)f^n(y)\,dy,
\]

we obtain

\begin{equation}
\tag{4.41}
v_n(x)=D_n\int_{B_2^+\cup B_2^-}\Gamma(x-y)f^n(y)\,dy.
\end{equation}

It now follows from (4.36) and (4.38), applied to (4.40) and (4.41),

\begin{equation}
\tag{4.42}
|Dv|_{0,\alpha;B_1^+}'\le C|f|_{0,\alpha;B_2^+}', \qquad C=C(n,\alpha).
\end{equation}

\end{theorem}

\begin{theorem}\label{gt:thm-poisson-first-derivative-halfball}\dcref{Theorem 4.16, (4.39), (4.40), (4.41), (4.42), (4.43), (4.44), (4.45), (4.46)}\group{ch04-poisson-newtonian-potential}\source{gilbarg-trudinger:page-0081, gilbarg-trudinger:page-0082}
\textit{Let $u\in C^0(\overline{B_2^+})$ satisfy Poisson's equation (4.33), where $f\in C^\alpha(\overline{B_2^+})$, and suppose $u=0$ on $B_2\cap\{x_n=0\}$. Then}

\begin{equation}
\tag{4.43}
|u|_{1,\alpha;B_1^+}'\le C\bigl(|u|_{0;B_2^+}+R|f|_{0,\alpha;B_2^+}'\bigr), \qquad C=C(n,\alpha).
\end{equation}

The proof is obtained from (4.42) by the concluding argument in Theorem 4.11.
It is also possible to state analogues of Theorems 4.3, 4.13 and Corollary 4.14 for solutions of (4.33). The details are left to the reader.
The preceding results can also be extended to equations of the form

\begin{equation}
\tag{4.44}
\Delta u=g+\operatorname{div}f
\end{equation}

when $f$ is Hölder continuous and $g$ is bounded and integrable. We observe that the first derivatives of the Newtonian potential of $g$ satisfy an $\alpha$-Hölder estimate for every $\alpha<1$ (see Problem 4.8(a); also Theorem 3.9). From this estimate for the

% PDF page 82
\source{gilbarg-trudinger:page-0082}

Newtonian potential, it follows that the $C^{1,\alpha}$ estimates for solutions of (4.44) corresponding to (4.37) and (4.43) take the form

\begin{equation}
\tag{4.45}
\|u\|_{1,\alpha;B_1} \leq C\bigl(\|u\|_{0;B_2} + R^2\|g\|_{0;B_2} + R\|f\|_{0,\alpha;B_2}\bigr)
\end{equation}

\begin{equation}
\tag{4.46}
\|u\|_{1,\alpha;B_1^+} \leq C\bigl(\|u\|_{0;B_2^+} + R^2\|g\|_{0;B_2^+} + R\|f\|_{0,\alpha;B_2^+}\bigr)
\end{equation}

where $C = C(n,\alpha)$.

\textit{Remark.} If $g \in L^p(\Omega)$, where $p = n/(1-\alpha)$, then the terms containing $g$ on the right-hand side in (4.45), (4.46) can be replaced by $R^{1+\alpha}\|g\|_p$ (see Problem 4.8(b)).
\end{theorem}

\section*{Notes}

The Hölder estimates of this chapter are essentially due to Korn [KR 1].

In Lemma 4.2, Hölder continuity can be replaced by Dini continuity, so that the Newtonian potential of $f$ is a $C^2$ solution of Poisson's equation $\Delta u = f$ if

\begin{equation}
\tag{4.47}
|f(x)-f(y)| \leq \varphi(|x-y|), \quad \text{where } \int_{0^+} \varphi(r)r^{-1}\,dr < \infty;
\end{equation}

(see Problem 4.2). However, if $f$ is only continuous the Newtonian potential need not be twice differentiable; (see Problem 4.9).

The weighted interior norms and seminorms (4.17), (4.29) are adapted from Douglas and Nirenberg [DN]. The primary function of the partially interior norms and seminorms (4.29) is to facilitate the derivation of boundary estimates by direct imitation of the proofs of interior estimates; (see, for example, Theorem 4.12 and Lemma 6.4).

\section*{Problems}

\textbf{4.1.} \ (a) Prove (4.7). (b) If $f \in C^{\alpha}(\mathbb{R})$ and $g \in C^{\beta}(\Omega)$, show that $f \circ g \in C^{\alpha\beta}(\Omega)$.

\textbf{4.2.} \ Prove Lemma 4.2 if $f$ is Dini continuous in $\Omega$ (i.e., $f$ satisfies (4.47)).

\textbf{4.3.} \ Show that Theorem 4.3 continues to hold if the boundedness of $f$ is replaced by $f \in L^p(\Omega)$ for some $p > n/2$; (see Lemma 7.12).

\textbf{4.4.} \ Derive Theorem 4.6 from Theorem 4.5 by applying the representation formula (2.17) to $u(x)=u(x)\eta(|x-x_0|/R)$, where $\eta$ is a cut-off function such that $\eta \in C^2(\mathbb{R})$, $\eta(r)=1$ for $r \leq \frac{3}{2}$, $\eta(r)=0$ for $r \geq 2$.

% PDF page 83
\source{gilbarg-trudinger:page-0083}

\textbf{4.5.} \ Prove the following extensions to Poisson's equation of the solid mean value inequalities (2.6) for Laplace's equation. Let \(u \in C^2(\Omega) \cap C^0(\bar\Omega)\) satisfy \(\Delta u = (\ge,\le) f\) in \(\Omega\). Then for any ball \(B = B_R(y) \subset \Omega\), we have
\[
u(y) = (\le,\ge)\left\{\frac{1}{\lvert B\rvert}\int_B u\,dx - \frac{1}{n\omega_n}\int_B f(x)\Theta(r,R)\,dx\right\},\qquad r = \lvert x-y\rvert,
\]
where
\[
\Theta(r,R)=
\begin{cases}
\dfrac{1}{n-2}\bigl(r^{2-n}-R^{2-n}\bigr)-\dfrac{R^2-r^2}{2R^n}, & n>2,\\[6pt]
\log(R/r)-\dfrac{1}{2}\bigl(1-r^2/R^2\bigr), & n=2.
\end{cases}
\]

\medskip

\noindent\textbf{4.6.} \ Prove Theorem 4.9 if the ball \(B\) is replaced by an arbitrary bounded \(C^2\) domain (use Lemma 14.16 and a comparison function \(d^m\eta\), where \(\eta\) is a suitable cut-off function and \(d\) is the distance function).

\medskip

\noindent\textbf{4.7.} \ Let \(\Delta u = f\) in \(\Omega \subset \mathbb{R}^n\). Show that the Kelvin transform of \(u\), defined by
\[
v(x)=\lvert x\rvert^{2-n}u(x/\lvert x\rvert^2)\qquad\text{for } x/\lvert x\rvert^2 \in \Omega
\]
satisfies
\[
\Delta v(x)=\lvert x\rvert^{-n-2}f(x/\lvert x\rvert^2).
\]

\medskip

\noindent\textbf{4.8.} \ Let \(w\) be the Newtonian potential of \(f\) in \(B = B_R(x_0)\).
\begin{enumerate}
\item[(a)] If \(f \in L^\infty(B)\), show that \(Dw \in C^\alpha(\mathbb{R}^n)\) for every \(\alpha \in (0,1)\) and
\[
[Dw]_{\alpha;B} \le C(n,\alpha)R^{1-\alpha}\lVert f\rVert_{\infty;B}
\]
\item[(b)] If \(f \in L^p(B)\), where \(p = n/(1-\alpha)\), \(0<\alpha<1\), show that \(Dw \in C^\alpha(\mathbb{R}^n)\) and
\[
[Dw]_{\alpha;B} \le C(n,\alpha)\lVert f\rVert_{p;B}.
\]
\end{enumerate}

\medskip

\noindent\textbf{4.9.}
\begin{enumerate}
\item[(a)] For \(\alpha\) with \(\lvert \alpha\rvert = 2\) let \(P\) be a homogeneous harmonic polynomial of degree \(2\) with \(D^\alpha P \neq 0\) (e.g. \(P = x_1x_2\), \(D_{12}P = 1\)). Choose \(\eta \in C_0^\infty(\{x\mid \lvert x\rvert < 2\})\) with \(\eta = 1\) when \(\lvert x\rvert < 1\), set \(t_k = 2^k\), and let \(c_k \to 0\) as \(k \to \infty\), with \(\sum c_k\) divergent. Define
\[
f(x)=\sum_{k=0}^\infty c_k \Delta(\eta P)(t_kx).
\]
Show that \(f\) is continuous but that \(\Delta u = f\) does not have a \(C^2\) solution in any neighbourhood of the origin.
\end{enumerate}

% PDF page 84
\source{gilbarg-trudinger:page-0084}

(b) For $\alpha$ with $|\alpha|=3$ choose a homogeneous harmonic polynomial $Q$ of degree $3$ with $D^\alpha Q\neq 0$. With $\eta$, $t_k$ and $c_k$ as in part (a), define

\[
u(x)=\sum_{0}^{\infty} c_k\eta(t_k x)Q(x)=\sum_{0}^{\infty} c_k(\eta Q)(t_kx)/t_k^3.
\]

Then

\[
\Delta u=g(x)=\sum_{0}^{\infty} c_k \Delta(\eta Q)(t_k x)/t_k.
\]

Show that $g\in C^1$ but that $u\notin C^{2,1}$ in any neighbourhood of the origin. Hence Lemma 4.4 is not valid for $\alpha=1$.

\medskip

\noindent\textbf{4.10.} Let $u\in C_0^2(B)$ satisfy $\Delta u=f$ in $B=B_R(x_0)$. Show that

\[
\text{(a)}\ |u|_0\le \frac{R^2}{2n}|f|_0;\qquad
\text{(b)}\ |D_i u|_0\le R|f|_0,\quad i=1,\ldots,n.
\]

Hence in (4.14), $|u|_{1;B}\le 3R^2|f|_{0;B}$.


% PDF page 85
\source{gilbarg-trudinger:page-0085}

Chapter 5

Banach and Hilbert Spaces

This chapter supplies the functional analytic material required for our study of
existence of solutions of linear elliptic equations in Chapters 6 and 8. This material
will be familiar to a reader already versed in basic functional analysis but we shall
assume some acquaintance with elementary linear algebra and the theory of metric
spaces. Unless otherwise indicated, all linear spaces used in this book are assumed
to be defined over the real number field. The theory of this chapter, however, carries
over almost unchanged if the real numbers are replaced by the complex numbers.
Let $\mathcal V$ be a linear space over $\mathbb R$. A norm on $\mathcal V$ is a mapping $p: \mathcal V \to \mathbb R$ (henceforth
we write $p(x)=\|x\|=\|x\|_{\mathcal V},\ x\in \mathcal V$) satisfying

\[
\text{(i)}\ \|x\|\ge 0\text{ for all }x\in \mathcal V,\ \|x\|=0\text{ if and only if }x=0;
\]
\[
\text{(ii)}\ \|\alpha x\|=|\alpha|\,\|x\|\text{ for all }\alpha\in \mathbb R,\ x\in \mathcal V;
\]
\[
\text{(iii)}\ \|x+y\|\le \|x\|+\|y\|\text{ for all }x,y\in \mathcal V\text{ (triangle inequality).}
\]

A linear space $\mathcal V$ equipped with a norm is called a normed linear space. A normed
linear space $\mathcal V$ is a metric space under the metric $\rho$ defined by

\[
\rho(x,y)=\|x-y\|,\qquad x,y\in \mathcal V.
\]

Consequently a sequence $\{x_n\}\subset \mathcal V$ converges to an element $x\in \mathcal V$ if $\|x_n-x\|\to 0$.
Also $\{x_n\}$ is a Cauchy sequence if $\|x_n-x_m\|\to 0$ as $m,n\to \infty$. If $\mathcal V$ is complete,
that is every Cauchy sequence converges, then $\mathcal V$ is called a Banach space.

\textit{Examples.} \ (i) Euclidean space $\mathbb R^n$ is a Banach space under the standard norm:

\[
\|x\|=\left(\sum_{i=1}^n x_i^2\right)^{1/2},\qquad x=(x_1,\ldots,x_n).
\]

(ii) For a bounded domain $\Omega\subset \mathbb R^n$, the H\"older spaces $C^{k,\alpha}(\bar\Omega)$ are Banach
spaces under either of the equivalent norms (4.6) or (4.6)$'$ introduced in Chapter 4;
(see Problems 5.1, 5.2).

(iii) The Sobolev spaces $W^{k,p}(\Omega)$, $W_0^{k,p}(\Omega)$ (see Chapter 7).

Existence theorems in partial differential equations are often reducible to the
solvability of equations in appropriate function spaces. For the Schauder theory of
linear elliptic equations we will employ two basic existence theorems for operator
